\chapter{The Likelihood}
\label{ch:loglikelihood}

\newthought{The likelihood closes the per-point loop.} After the workflow has filled the
observables dictionary, a list of named expressions turns it into a score. This chapter
covers where those expressions live, how the total \code{LogL} is formed, and what happens
when evaluation fails.

\section{Where likelihood terms live}
\label{sec:likelihood-blocks}

Two equivalent homes, one precedence rule:

\begin{yamlwin}
Likelihood:                      # preferred, top level
  expressions:
    - name: LogL_Z
      expression: LogGauss(z, 1.0, 0.1)
    - name: LogL_mass
      expression: -0.5 * ((mh - 125.25) / 0.17)**2
\end{yamlwin}

\begin{yamlwin}
Sampling:                        # accepted alias, lower precedence
  LogLikelihood:
    - name: LogL_Z
      expression: z
\end{yamlwin}

\yamlkey{Likelihood.expressions} wins when both are present --- first non-empty list
wins; they are \emph{never merged}. Each entry:

\begin{keylist}
  \item[name] \opt, default \code{LogL}. The term's observable name; results are added
        to the observables dictionary and appear in the output records.
  \item[expression] \req. An expression in the shared language of
        Chapter~\ref{ch:expressions}, evaluated over the observables.
\end{keylist}

\section{How \code{LogL} is formed}
\label{sec:logl-semantics}

Terms evaluate \emph{in order}, and each term's result is added to the observables before
the next term runs --- so later terms can reference earlier ones by name. Then:

\begin{itemize}
  \item If some term is literally named \code{LogL}, that term \textbf{is} the total ---
        no summing. Use this when you want full control:
\begin{yamlwin}
Likelihood:
  expressions:
    - name: LogL_a
      expression: LogGauss(obs_a, 5.1, 0.4)
    - name: LogL_b
      expression: LogGauss(obs_b, 0.7, 0.05)
    - name: LogL
      expression: Max(LogL_a, LogL_b)   # explicit total
\end{yamlwin}
  \item Otherwise, $\code{LogL} = \sum_i \text{term}_i$ --- the ordinary case where
        independent measurement terms add.
\end{itemize}

When neither likelihood block exists, the execution plan simply carries no likelihood
step and \code{LogL} is absent from the records --- legitimate for pure mapping scans
where you only want observables.

\begin{jnote}
Everything is \emph{log}-likelihood, natural log. The convenience kernels differ in their
normalisation: \code{LogGauss(x, mu, sigma)} is
$-\tfrac{1}{2}\big(\tfrac{x-\mu}{\sigma}\big)^2$ without the constant term;
\code{Normal} includes the $1/(\sigma\sqrt{2\pi})$ factor; \code{Gauss} is the
unnormalised kernel $e^{-\frac{1}{2}((x-\mu)/\sigma)^2}$ (not a log!). See
Chapter~\ref{ch:expressions}.
\end{jnote}

\section{Failure semantics}
\label{sec:likelihood-failure}

A likelihood expression that references an observable missing from the dictionary raises
a structured \emph{missing-variables} error naming the absent symbols, and the Sample
\textbf{fails} --- it is recorded as failed, keeps its logs, and the scan continues.
This is deliberate: a missing observable usually means an upstream calculator failed to
produce output for this point, and silently scoring the point would poison the scan.

If a term should be optional --- an observable only some points produce --- make the
producing module emit an explicit default instead of omitting the key
(Chapters~\ref{ch:calculators} and~\ref{ch:operas}: absent output entries are captured
as null values rather than dropped, so the expression sees \emph{something}).

\section{Who consumes \code{LogL}}
\label{sec:logl-consumers}

For the fixed-set methods, the likelihood is bookkeeping: every accepted point runs
regardless, and \code{LogL} lands in the database for your posterior analysis. Two
consumers act on it at runtime:

\begin{itemize}
  \item \sampler{AdaptiveBridson} evaluates its \yamlkey{target\_expression} over the
        returned observables --- including likelihood terms --- to steer refinement
        (Chapter~\ref{ch:adaptive-levelset}).
  \item Your plots: \file{samples.csv} carries every term as a column, so per-term
        pulls are one \code{pandas} expression away.
\end{itemize}

\begin{jtip}
Name terms by what they constrain --- \code{LogL\_mh}, \code{LogL\_relic},
\code{LogL\_bsg} --- and let the automatic sum form the total. Per-term columns in the
output are the fastest way to answer ``which constraint kills this region?''.
\end{jtip}
